% titel: Blatt 3, A5: Gleichgewicht auf der schiefen Ebene
% skript: Kap. Statik, Haft- und Gleitreibung

% Haftreibung, Gleitreibung, Reibungswinkel.

Auf einer geneigten Ebene halten sich Gewicht, Normalkraft und Reibung
das Gleichgewicht; die Grenzwoerte der Reibung bestimmen den Winkel, bei
dem der Koerper rutscht.

\begin{werkzeugkasten}[Werkzeuge]
\textbf{Komponenten der Gewichtskraft} ($\vec F_G$ zerlegt senkrecht und
parallel zur Ebene):
\begin{equation}
  F_\parallel=mg\sin\alpha,\qquad
  F_\perp=mg\cos\alpha=F_N.
\end{equation}

\textbf{Haftreibung} (Solange der Koerper ruht):
\begin{equation}
  |\vec F_R|\le F_{R,\max}=\mu_H F_N.
\end{equation}
Grenzwert beim Anlaufen des Rutschens (Haftreibungswinkel):
\begin{equation}
  \mu_H=\tan\alpha_H\;.
\end{equation}

\textbf{Gleitreibung} (gleitender Koerper):
\begin{equation}
  |\vec F_R|=\mu_G F_N,\qquad\mu_G=\tan\alpha_G
\end{equation}
im gleichfoermigen Gleiten ($a=0$).

\textbf{Beschleunigung beim Rutschen} (Antrieb minus Gleitreibung als
Massenkraft):
\begin{equation}
  a=g(\sin\alpha-\mu_G\cos\alpha)\;.
\end{equation}
\end{werkzeugkasten}

\begin{werkzeugkasten}[Vorgehen]
\begin{itemize}
  \item Freikoerperbild: $\vec F_G$, $\vec F_N$ senkrecht, $\vec F_R$ den
        Hang hinauf, solange der Koerper (noch) ruht.
  \item Haftreibung: nur der \emph{maximale} Wert $\mu_H F_N$ haelt an;
        Gleichgewicht bis $\tan\alpha=\mu_H$.
  \item Gleitreibung: konstanter Wert $\mu_GF_N$; $a=g(\sin\alpha-\mu_G\cos\alpha)$.
  \item Eine Strasse darf nur bis zum Winkel mit $\tan\alpha_{\max}=\mu_H$
        mit blockierten Raedern geneigt sein (kein Rutschen).
\end{itemize}
\end{werkzeugkasten}

\begin{achtung}[Typische Fallen]
\begin{itemize}
  \item $\mu_H\ge\mu_G$; ein Rutschen mit $\mu_G>\mu_H$ ist unphysikalisch
        (s. unten).
  \item $F_N=mg\cos\alpha$, nicht $mg$; bei steiler Ebene nimmt die
        Normalkraft ab.
  \item Die Masse kuerzt sich in $a=g(\sin\alpha-\mu\cos\alpha)$ heraus;
        die Auflageflaeche spielt im einfachen Reibungsmodell keine Rolle.
  \item ``Blockiert und parkiert'' bedeutet Haftreibung ($\mu_H$), nicht
        $\mu_G$.
  \item Die Hoehe der Haftreibung am Berg ohne Ansto\ss{} ist nur bis zur
        Grenze $\mu_HF_N$ bestimmt; steiler wird sie nicht.
\end{itemize}
\end{achtung}
